%%

%

%%%%

%%%   Самарский государственный технический университет

%%%   Кафедра прикладной математики и информатики

%%%

%%%   Преамбула для статьи в журнал "Вестник Самарского государственного

%%%   технического университета. Серия: «Физико-математические науки»"

%%%   Ориентирована под MikTEx 2.4 for Win32

%%%

%%%   Хотя сам журнал собирается под GNU/Linux на дистрибутиве TeXLive



\documentclass[11pt,a4paper,twoside]{article}



\usepackage[cp1251]{inputenc}

\usepackage[T2A]{fontenc}

\usepackage[english,russian]{babel}

\usepackage{samgtubib}

\usepackage[dvips]{graphicx}

\usepackage[multiple]{footmisc}



\usepackage{samgtu}

\renewcommand{\VestnikYear}{2009} % Год издания Вестника

\renewcommand{\VestnikNumber}{2\,(19)} % Номер Вестника

\renewcommand{\VestnikMonth}{Ноябрь} % Месяц выхода Вестника

\renewcommand{\VestnikData}{01.11.09} % Дата подписания Вестника в печать

\renewcommand{\VestnikUPL}{26,64} % Усл. печ. л.

\renewcommand{\VestnikUIL}{25,73} % Уч.-изд. л.

\renewcommand{\VestnikRegNumber}{479/09} % Регистрационный номер

\renewcommand{\VestnikOrder}{994} % Номер заказа







%

%%%%%

%%%   Самарский государственный технический университет

%%%   Кафедра прикладной математики и информатики

%%%

%%%   Преамбула для статьи в журнал "Вестник Самарского государственного

%%%   технического университета. Серия: «Физико-математические науки»"

%%%   Ориентирована под MikTEx 2.4 for Win32

%%%

%%%   Хотя сам журнал собирается под GNU/Linux на дистрибутиве TeXLive



\documentclass[11pt,a4paper,twoside]{article}



\usepackage[cp1251]{inputenc}

\usepackage[T2A]{fontenc}

\usepackage[english,russian]{babel}

\usepackage{samgtubib}

\usepackage[dvips]{graphicx}

\usepackage[multiple]{footmisc}



\usepackage{samgtu}

\renewcommand{\VestnikYear}{2009} % Год издания Вестника

\renewcommand{\VestnikNumber}{2\,(19)} % Номер Вестника

\renewcommand{\VestnikMonth}{Ноябрь} % Месяц выхода Вестника

\renewcommand{\VestnikData}{01.11.09} % Дата подписания Вестника в печать

\renewcommand{\VestnikUPL}{26,64} % Усл. печ. л.

\renewcommand{\VestnikUIL}{25,73} % Уч.-изд. л.

\renewcommand{\VestnikRegNumber}{479/09} % Регистрационный номер

\renewcommand{\VestnikOrder}{994} % Номер заказа







%

%\usepackage{qrcode}

%

%%

%%\newcommand{\totop}[1]{\raisebox{6pt}[1pt][2pt]{#1}} 

%%\newcommand{\tottop}[1]{\raisebox{12pt}[1pt][2pt]{#1}} 

%%\newcommand{\totttop}[1]{\raisebox{18pt}[1pt][2pt]{#1}} 

%%\newcommand{\tottttop}[1]{\raisebox{24pt}[1pt][2pt]{#1}} 

%%\newcommand{\totttttop}[1]{\raisebox{30pt}[1pt][2pt]{#1}} 

%\newcommand{\tobot}[1]{\raisebox{-6pt}[1pt][2pt]{#1}} 

%%\newcommand{\tobbot}[1]{\raisebox{-12pt}[1pt][2pt]{#1}} 

%%\newcommand{\tobbbot}[1]{\raisebox{-18pt}[1pt][2pt]{#1}}

%

%\graphicspath{{./00PIC/}}

%

%%\samgtupubl % для генерации pdf, отдаваемого в типографию СамГТУ

%\pdfcrop % обрезка pdf для размещения на math-net.ru

%\draft % На корректуре

%\filllastpage % Последняя страница не заполнена не полностью

%%\briefcomm % Статья является кратким сообщением



\vsgtu{1491}



\FirstPage{410}

\LastPage{422}

%%

%%

%\begin{document}





\renewcommand{\baselinestretch}{0.9}



\hfill \qrcode[hyperlink,height=.8in]{http://dx.doi.org/10.14498/vsgtu1491}

\vspace{-.8in}



%%%%%%%%%%%%%%%%%%%%%%%%%%%%%%%%%%%%%%%%%%%%%%%%%%%%%%%%%%%%%%%%%%%%%%%%%%%%%%%%%%%%

%Информация о статье и об авторах на русском и английском языках



%%%%%%%%%%%%%%%%%%%%%%%%%%%%%%%%%%%%%%%%%%%%%%%%%%%%%%%%%%%%%%%%%%%%%%%%%%%%%%%%%%%%

%Информация о статье и об авторах на русском и английском языках



% Номер УДК

\udc{517.956.2}%Обработка текста. Подготовка текстов : Языки программирования. Компьютерные языки

\msc{35J25, 35J56}% Коды MSC см. http://www.ams.org/msc/



\rutitle{О решениях задачи  Шварца для $J$-аналитических функций с~одинаковым жордановым базисом действительной и мнимой части матрицы~$J$}% Полный заголовок

{О решениях задачи  Шварца для $J$-аналитических функций\dots}% Заголовок, который идёт в верхний колонтитул



\titlehack{О решениях задачи  Шварца для~$\boldsymbol  J$-аналитических~функций с~одинаковым жордановым базисом действительной и~мнимой~части матрицы $\boldsymbol J$}



\ruauthor{В.~Г.~Николаев}{Н\,и\,к\,о\,л\,а\,е\,в~В.~Г.}



% Если несколько, то так  \institute{ Организация 1 \and Организация 2  \and Организация 3}



%\email{E-mail}{vg14@inbox.ru} %E-mail's авторов



\ruaffil{\runovgu.}

\enaffil{\ennovgu.}





\ruauthordetails{1}{\fullauthor{\rupid{78216}{Владимир Геннадьевич Николаев}}\regalia{к.ф.-м.н.; \mem{vg14@inbox.ru}}

\position{доцент} \dept{каф. алгебры и геометрии}}



\enauthordetails{1}{\fullauthor{\enpid{78216}{Vladimir G. Nikolaev}} \regalia{Cand. Phys. \& Math. Sci.; \mem{vg14@inbox.ru}} \position{Associate Professor}

\dept{Dept. of Algebra and Geometry}}



\rukeywords{задача Шварца, $\lambda$-голоморфная  функция, контур Ляпунова,  жорданова форма, жорданов базис,

характеристическое уравнение, эллиптическая система, задача Дирихле}



\ruabstract{Изучена граничная задача Шварца для $J$-аналитических функций.

Они представляют собой решения линейной  комплексной системы дифференциальных уравнений   в~частных производных первого порядка.

Рассмотрен   тот случай, когда действительная и мнимая части матрицы $J$ приводятся к треугольному виду одним и~тем же комплексным преобразованием.

В основной теореме доказан  критерий для  собственных чисел матрицы $J$, при выполнении которого в~плоских областях, ограниченных контуром Ляпунова, решение задачи Шварца существует и~единственно.

Так же получена равносильная форма данного критерия,

которая во многих случаях  более удобна для проверки.

При доказательстве  теоремы используются известные факты о граничных свойствах $\lambda$"=голоморфных функций.

Само доказательство основано на методе прямой и~обратной

редукции  задачи Шварца к~задаче

   Дирихле для  вещественных эллиптических систем в частных производных

   второго порядка.

Построены  примеры  матриц, для которых указанный критерий  выполнен.}

%

\entitle{On decisions of Schwartz' problem for $J$-analytic functions with the same Jordan basis of real and imaginary parts of $J$-matrix}



\engtitlehack{On decisions of Schwartz' problem for $\boldsymbol J$-analytic functions with the same Jordan basis of real and~imaginary parts of $\boldsymbol J$-matrix}





\enauthor{V.~G.~Nikolaev}{N\,i\,k\,o\,l\,a\,e\,v~V.~G.}





\enabstract{Boundary Schwartz' problem for $J$-analytic functions was studied within this scientific work. These functions are solutions of linear complex system of partial differential equations of the first order. It was considered, that the real and imaginary parts of $J$-matrix are put into triangular form by means of one and the same complex transformation. The main theorem proved a criterion for eigenvalues of $J$-matrix. Shall this criterion be fulfilled within the complex plane within the boundaries defined by Lyapunov line, there is a decision on Schwartz' problem and it is the only one. The equal form of this criterion was found, which in many cases is more convenient for check. While proving the theorem, known facts about boundary properties of $l$-holomorphic functions are applied. The proof itself is based on the method of direct and reverse reduction of Schwarz' problem to Dirichlet's problem for real valued elliptic systems of partial differential equations of the second order. Examples of matrices are given, whereby the specified criterion is fulfilled.}



\enkeywords{Schwartz' problem, $l$-holomorphic function, Lyapunov line, Jordan form, Jordan basis, characteristic equation, elliptic system, Dirichlet's problem}



\date{14/IV/2016} % Дата поступления статьи в редакцию.

\revisiondate{22/V/2016}

\accepteddate{27/V/2016}



%%%%%%%%%%%%%%%%%%%%%%%%%%%%%%%%%%%%%%%%%%%%%%%%%%%%%%%%%%%%%%%%%%%%%%%%%%%%%%%%%%%%





\makerutitle



\Section[n]{Введение}

Обобщением  понятия голоморфных функций являются функции,

 аналитические по Дуглису (A.~Douglis) [\citen{1Nk-S1}--\citen{4Nk-S3}].



\smallskip



\hypertarget{nik:def1}{}



\begin{definition}[1~\cite{1Nk-S1}]Пусть  матрица $J\in\mathbb{C}^{n\times n}$ не имеет вещественных собственных чисел$.$ 

Аналитической по Дуглису $($или $J$-аналитической с матрицей $J)$ называется   комплексная $n$"=вектор"=функция

 $ \phi= \phi(z)\in C^1(D),$ для которой в~области  $ D\subset\mathbb{R}^2$  выполнено  уравнение

\begin{equation}

\frac{\partial\phi}{\partial y}-J\cdot\frac{\partial\phi}{\partial x}=0,\quad

z\in D.

\label{1n}

\end{equation}

\end{definition}



\smallskip



%\begin{theorem}[Виета]

%Если $x_1$ и~$x_2$ "--- корни квадратного уравнения $$x^2+px+q,$$ то $x_1 + x_2 = -p$ и~$x_1 x_2 = q.$

%\end{theorem}

%

%\begin{proof} Для доказательства достаточно рассмотреть произведение $(x-x_1)(x-x_2)$.

%\end{proof}

%

%

%\begin{theorem}

%Для любых $x$ и~$y$ из~$\mathbb R$ справедливо $(x+y)^2 = x^2+2xy+y^2$.

%\end{theorem}

%

%\begin{theorem}[1]

%Для любых $x$ и~$y$ из~$\mathbb R$ справедливо $(x+y)^2 = x^2+2xy+y^2$.

%\end{theorem}

%

%\begin{theorem}[1, Н.И. Мусхелишвили]

%Для любых $x$ и~$y$ из~$\mathbb R$ справедливо $(x+y)^2 = x^2+2xy+y^2$.

%\end{theorem}

%

%\begin{lemma}[1]

%Для любых $x$ и~$y$ из~$\mathbb R$ справедливо $(x+y)^2 = x^2+2xy+y^2$.

%\end{lemma}

%

%\begin{lemma}[5]

%Для любых $x$ и~$y$ из~$\mathbb R$ справедливо $(x+y)^2 = x^2+2xy+y^2$.

%\end{lemma}

%

%\begin{example}[1]

%Для любых $x$ и~$y$ из~$\mathbb R$ справедливо $(x+y)^2 = x^2+2xy+y^2$.

%\end{example}

%

%\begin{example}[3]

%Для любых $x$ и~$y$ из~$\mathbb R$ справедливо $(x+y)^2 = x^2+2xy+y^2$.

%\end{example}

%

%\begin{definition}[1]

%gjvadjkvlerk

%\end{definition}

%

%\begin{definition}[2]

%gjvadjkvlerk

%\end{definition}



%Аналогично введены окружeния \verb"remark" для замечаний; \verb"example" для примеров, \verb"lemma" для лемм, \verb"definition" для определений.



%\end{document}





\hypertarget{nik:def2}{}

\begin{definition}[2]В~скалярном случае при $ J=\lambda,$

${\rm Im}\,\lambda\ne 0~$ функцию

$ f_\lambda(z)\in C^1(D),$ для которой в   области  $ D\subset\mathbb{R}^2$

выполнено уравнение

\begin{equation}

\frac{\partial f_\lambda}{\partial y}-\lambda\cdot

\frac{\partial f_\lambda}{\partial x}=0,\quad z\in D,

\label{2n}

\end{equation}

  назовем      $ \lambda$-голоморфной в области $ D.$

  \end{definition}





\smallskip



При $\lambda=i~$  функции \eqref{2n} совпадают с обычными голоморфными.



\smallskip



\begin{definition}[3] Будем говорить, что функция $ \phi(z)$ соответствует матрице $ J,$ если  выполнено  уравнение  \eqref{1n}.

\end{definition}



\smallskip



Как показано в  \cite{2Nk-Nk2}, комплексная система дифференциальных

уравнений в~частных производных первого порядка \eqref{1n} является эллиптической.

 Рассмотрим для системы \eqref{1n}   следующую граничную  задачу Шварца

  [\citen{1Nk-S1,2Nk-Nk2}].



\smallskip

{\small\sc Задача Шварца.}

 {\it Пусть конечная односвязная область $ D\subset\mathbb{R}^2$

 ограничена  гладким контуром $ \Gamma.$ Требуется

найти $J$-аналитическую с матрицей $J$

 в~области $D$ функцию

 $\phi(z)=u(x,y)+iv(x,y)\in C(\overline{D}),$ которая

 удовлетворяет краевому  условию

\begin{equation}

{\rm Re}\,\phi(z)\bigr|_{\Gamma}=\varphi(z),

\label{3n}

\end{equation}

 где вещественная граничная вектор-функция

 $\varphi(z)=\bigr(\varphi_1,\ldots,\varphi_n\bigr)\in C(\Gamma)$ задана}.



\smallskip



При $\varphi\equiv 0$ будем говорить об {\it однородной  задаче} Шварца.

Очевидными решениями последней служат постоянные векторы

$\phi=i\overline{c},~$ $\overline{c}\in \mathbb{R}^{n},$

 которые принято называть  {\it тривиальными решениями.}



Отметим, что   задача Шварца  \eqref{3n}

равносильна краевым задачам для многих видов

эллиптических систем

дифференциальных уравнений в частных производных первого и второго порядков

[\citen{4Nk-S3}--\citen{9Nk-B3}].

 Эти системы, в свою очередь,  имеют широкое применение в теории упругости

 [\citen{10Nk-Bo-1}--\citen{13Nk-J2}], поэтому

изучение задачи Шварца имеет как чисто теоретический, так и прикладной интерес.





Как известно, при $n=1$ {\it однородная }

задача \eqref{3n} имеет только тривиальные решения при любом типе границы

$\Gamma$ области $D.$

{\it Неоднородную } задачу Шварца удобно рассматривать \cite{2Nk-Nk2}

в областях, ограниченных контуром Ляпунова.

Поэтому приведем следующее известное



\smallskip

\begin{definition}[4~\cite{14Nk-Mus}]

 Гладкая кривая $\Gamma\subset\mathbb{R}^2$ называется кривой Ляпунова, если

    существуют  такие два вещественных числа $a>0$ и $b,$ $0<b\le 1,$

   что для любых двух точек $z_1,z_2\in \Gamma$   выполняется

    условие Ляпунова

    $$

    |\theta|<a\cdot|z_1-z_2|^b,

      $$

     где $\theta$  "--- угол между касательными или нормалями

      к $\Gamma$ в  точках $z_1, z_2.$

      \end{definition}

      

\smallskip



Примерами  кривых Ляпунова могут служить окружность и эллипс.



При  $n=1$  имеет место следующая теорема существования и единственности

решения неоднородной задачи Шварца.



\smallskip



\hypertarget{nik:th1}{}



 \begin{theorem}[1 (Н.\,И.~Мусхелишвили, \cite{14Nk-Mus})]Пусть $\Gamma=\partial D$ 

 "--- контур Ля\-пу\-но\-ва$,$

  граничная функция $\varphi(z)$ принадлежит  классу

Гельдера $H^\sigma(\Gamma),$ \mbox{$0<\sigma<1.$}

 Тогда

  можно построить единственную $($с точностью до постоянной$)$

 голоморфную в $D$ функцию $f(z)\in H^\sigma(\overline{D})$

  по  условию ${\rm Re}\,f(z)\bigr|_\Gamma=\varphi(z).$

\end{theorem}



\smallskip



 Обозначим $\lambda=a+bi,$ $a$, $b\in\mathbb{R}$, $b\ne 0.$ 

 Как известно,

 голоморфная  функция $f(x,y)$ становится $\lambda$-голоморфной

 $f_\lambda(x',y')$

  в результате  линейного  обратимого преобразования

  $x=x'+ay'$, $y=by'.$

  Поэтому теорема~\hyperlink{nik:th1}{1} справедлива и~для $\lambda$"=голоморфных

  функций $f_\lambda(z).$





Однако  при $n\ge 2$ вопрос существования  и единственности

решений задачи Шварца

уже намного сложнее. В частности, известны примеры нетривиальных

решений однородной задачи Шварца \cite{15Nk-Nk2}.

Приведем один из них.



\smallskip



\begin{example}[1]Пусть

$$

J=\left(\begin{array}{lr}

1-2i& 48-8i \\

\frac{\textstyle 1}{\textstyle 2}&2+8i

\end{array}\right),

$$

$$

\phi(z)=

\left(\begin{array}{r}

12x^2+32xy+46y^2-1+(-x^2+22xy+27y^2)\cdot i\\

(3x^2+8xy+30y^2)\cdot\frac{\textstyle i}{\textstyle 4}

\end{array}\right).

$$





Функция $\phi(z)$ соответствует матрице $J,$ которая

        имеет  разные

        собственные числа $\lambda=2+4i,~$ $\mu=1+2i.$

Как нетрудно видеть, ${\rm Re}\,\phi(z)\bigr|_{\Gamma}=0$

       на эллипсе $\Gamma$: $12x^2+32xy+46y^2=1.$

\end{example}



\smallskip



\Section{Редукция задачи Шварца к задаче Дирихле для эллиптических систем}

Ниже (теорема~\hyperlink{nik:th4}{4}) будет доказан аналог теоремы~\hyperlink{nik:th1}{1}  для  вектор"=функций,

аналитических по Дуглису со специальной матрицей $J$ при $n\ge 2.$ 

Прежде чем сформулировать данную теорему, выполним следующее преобразование задачи Шварца.





Обозначим $\phi=u+iv,$ где $u=\mathop{\mathrm{Re}}\phi$ и $v=\mathop{\mathrm{Im}}\phi$ ---

 вещественные $n$-вектор-функции,  и подставим матрицу

 \begin{equation}

J=A+Bi,\quad A,B\in\mathbb{R}^{n\times n},\quad \det B\ne 0

\label{4n}

\end{equation}

  в уравнение  \eqref{1n}.  

  Приравнивая к~нулю действительные и мнимые части, получим

\begin{equation}

     \frac{\partial u}{\partial y}-A \,\frac{\partial u}{\partial x}

+B\, \frac{\partial v}{\partial x}=0,\quad

\frac{\partial v}{\partial y}-A\, \frac{\partial v}{\partial x}

-B\, \frac{\partial u}{\partial x}=0.

\label{5n}

\end{equation}





Из  системы \eqref{5n}  выразим частные производные функции $v$:

\begin{equation}

\frac{\partial v}{\partial x}=B^{-1}A\,\frac{\partial u}{\partial x}

-B^{-1}\,\frac{\partial u}{\partial y},\quad

\frac{\partial v}{\partial y}=\bigr(B+AB^{-1}A\bigr)\,

\frac{\partial u}{\partial x}-

AB^{-1}\,\frac{\partial u}{\partial y}.

\label{6n}

\end{equation}



Из \eqref{6n} с учетом условия замкнутости

имеем

$$

\frac{\partial }{\partial y}\left[B^{-1}A\,\frac{\partial u}{\partial x}

-B^{-1}\,\frac{\partial u}{\partial y}\,\right]=

\frac{\partial }{\partial x}\left[\bigr(B+AB^{-1}A\bigr)\,

\frac{\partial u}{\partial x}-

AB^{-1}\,\frac{\partial u}{\partial y}\,\right],

$$

или

 \begin{equation}

 B\cdot(B+AB^{-1}A)\,\frac{\partial^2 u}{\partial x^2}-B\cdot

 (AB^{-1}+B^{-1}A)\,\frac{\partial^2 u}{\partial x\partial y}

 +\frac{\partial^2 u}{\partial y^2}=0.

 \label{7n}

\end{equation}



В работе \cite{2Nk-Nk2} показано, что   вещественная система \eqref{7n}

будет эллиптической.

Как нетрудно видеть, имеет место



\smallskip





\hypertarget{nik:lemma1}{}



\begin{lemma}[1]Если  существует решение

 $\phi=u+iv$ задачи Шварца $\eqref{3n},$ соответствующее

 данной матрице $J=A+Bi,$

то функция $u(x,y)$ будет решением  задачи Дирихле

\begin{equation}

u(x,y)\bigr|_\Gamma=\varphi(z)

\label{8n}

\end{equation}

для системы $\eqref{7n}.$

\end{lemma}



\smallskip



В обратную сторону утверждение леммы~\hyperlink{nik:lemma1}{1}  не всегда верно.

Но в частных случаях приведенная выше редукция   уравнения \eqref{1n}

к уравнению \eqref{7n}

обратима "--- на этом будет основано доказательство

основой   теоремы~\hyperlink{nik:th4}{4}.



\smallskip



\Section{Основная теорема}

 Сначала приведем формулировки двух теорем, которые будут   использованы ниже,    а также  сделаем  вспомогательные построения.



\smallskip



\hypertarget{nik:th2}{}

\begin{theorem}[2 (А.\,П. Солдатов, \cite{2Nk-Nk2})]Пусть  $\Gamma=\partial D$ "--- кривая Ляпунова и пусть

$(\mathop{\rm Im} \lambda)\cdot(\mathop{\rm Im} \mu)<0.$

  Тогда

 для любой граничной функции $\varphi(z)$ из класса

Гельдера $H^\sigma(\Gamma),$ $0<\sigma<1,$  решение задачи

 $$

f_\lambda(z)+f_\mu(z)=\varphi(z), \quad z\in\Gamma

$$

в классе  функции $f_\lambda(z),$ $f_\mu(z)\in H^\sigma(\overline{D})$

существует и  единственно $($с точностью до константы$).$

\end{theorem}



\smallskip



\hypertarget{nik:th3}{}



 \begin{theorem}[3 (Е.\,Ю. Панов, \cite{16Nk-Nk1})]Пусть $\Gamma=\partial D$ "--- жорданова  кривая$.$

 Пусть  $\lambda$- и $\mu$-голоморфные функции $f_\lambda(z), $ $f_\mu(z)\in{\rm Lip}(\overline{D}),$

 причем

 $ (\mathop{\rm Im}\lambda)\times(\mathop{\rm Im}\mu)<0$

 и $f_\lambda(z)=f_\mu(z)$ на $\Gamma.$ 

 Тогда

 $f_\lambda(z)=f_\mu(z)\equiv{\rm const}.$

 \end{theorem}



\smallskip

Предположим далее, что вещественные

матрицы $A$, $B$  в \eqref{4n} обладают  следующим свойством:

для некоторой матрицы

 $Q\in\mathbb{C}^{n\times n}$ матрицы

 $J_1=Q^{-1}AQ$ и $J_2=Q^{-1}BQ$  --- {\it нижне-треугольные. }

 Матрица $Q$ в частном случае может быть  жордановым базисом для

 $A$, $B$ (но может им и не быть). 

 Приведем наиболее простой

  пример таких матриц: $J=\alpha E+Ai$ и $J=A+\alpha iE,$

 где $\alpha\in\mathbb{R},$ $A\in\mathbb{R}^{n\times n}.$





Пусть $\overline{J}=A-Bi$ "--- сопряжение матрицы $J=A+Bi$

по аналогии с~сопряжением комплексного числа. Очевидно, что матрицы

\begin{equation}

\begin{array}{l}

J'=Q^{-1}JQ=Q^{-1}AQ+Q^{-1}BQi=J_1+J_2i,\\

J''=Q^{-1}\overline{J}Q=Q^{-1}AQ-Q^{-1}BQi=J_1-J_2i

\end{array}

\label{9n}

\end{equation}

тоже

будут  нижне-треугольными.



  Обозначим через   $\lambda_{1k}$, $\lambda_{2k}$ собственные  числа

  матриц $Q^{-1}JQ$,  $Q^{-1}\overline{J}Q$,

  соответственно,  стоящие на $k$-той строке матриц $J'$, $J''$ в~\eqref{9n}.





    Пусть  выполнены следующие неравенства:

\begin{equation}

\left(\mathop{\rm Im}\lambda_{1k} \right)\cdot \left(\mathop{\rm Im} \lambda_{2k} \right)<0,\quad  k=1,\ldots,n.

 \label{10n}

\end{equation}



Прежде чем сформулировать теорему~\hyperlink{nik:th4}{4},

преобразуем систему неравенств \eqref{10n} к равносильной форме \eqref{11n}.

Она не будет использована при доказательстве теоремы~\hyperlink{nik:th4}{4}, но удобна для практических

вычислений.



Обозначим через  $\lambda_k$, $\mu_k$,  $k=1,\ldots, n,$ собственные  числа   матриц $J_1$, $J_2$ в~\eqref{9n}

  соответственно (то есть собственные числа матриц $A$ и $B$),  стоящие на $k$-той строке. 

  Тогда с учетом \eqref{10n} имеем:

\begin{multline}

0>\left(\mathop{\rm Im} \lambda_{1k}\right)\cdot \left(\mathop{\rm Im} \lambda_{2k} \right)

=\mathop{\rm Im}\left(\lambda_k+i\mu_k \right)\cdot \mathop{\rm Im}\left(\lambda_k-i\mu_k\right)=

\\

=

\left(\mathop{\rm Im} \lambda_k+\mathop{\rm Re} \mu_k \right)\cdot

\left(\mathop{\rm Im} \lambda_k-\mathop{\rm Re} \mu_k \right)=

\left(\mathop{\rm Im} \lambda_k \right)^2-

\left(\mathop{\rm Re} \mu_k \right)^2,

\end{multline}

что равносильно системе неравенств

\begin{equation}

\left|\mathop{\rm Im} \lambda_k\right|<

\left|\mathop{\rm Re} \mu_k\right|,\quad

\quad  k=1,\ldots,n.

 \label{11n}

\end{equation}



Покажем, что справедлива



\smallskip



\hypertarget{nik:th4}{}

\begin{theorem}[4]Пусть матрица

$J=A+Bi,$ $A,$ $B\in\mathbb{R}^{n\times n},$

причем матрицы   $J_1=Q^{-1}AQ$ и $J_2=Q^{-1}BQ$ "---  нижне-треугольные

в некотором базисе $Q.$ 

Пусть также  выполнены неравенства  $\eqref{10n},$ или, что то же,

  $\eqref{11n}.$  

  Тогда

\begin{itemize}



\hypertarget{nik:th4-1}{}

\item[$1)$] если  жордановы  формы матриц $A,$ $B$ диагональны$,$ $Q$ "--- их жорданов базис

   и

 $\Gamma=\partial D$ "--- контур Ляпунова$,$ то

 для любой граничной вектор"=функции $\varphi(z)\in

 H^\sigma(\Gamma),$ $0<\sigma<1$  решение неоднородной задачи Шварца $\eqref{3n}$

 в классе  функций $\phi(z) \in H^\sigma(\overline{D})$

существует и  единственно $($с~точностью до вектор-константы$);$



\hypertarget{nik:th4-2}{}



\item[$2)$]  если  жордановы формы  матриц $A,$ $B$ произвольны  и

   $\Gamma=\partial D$ "--- жорданова кривая$,$

то решения однородной $(\varphi\equiv 0)$ задачи  Шварца $\eqref{3n}$

 в~классе  функций $\phi(z) \in {\rm Lip} (\overline{D})$ "---

только тривиальные$,$ то есть

$\phi\equiv i\overline{c},$ $\overline{c}\in \mathbb{R}^{n}.$



\end{itemize}

\end{theorem}



\smallskip



Таким образом, неравенства \eqref{10n}, \eqref{11n} можно считать

критерием (достаточным условием) существования и

единственности  решений задачи Шварца для функций, $J$-аналитических

с данным типом  матриц.



\smallskip



\begin{proof}  Покажем сначала, что из \eqref{10n}

вытекает условие

 $\det B\ne 0,$ которое необходимо для применения формулы \eqref{7n}.



\smallskip



 $\triangleright$ <<От противного>>: пусть $\det B=0$ при выполнении \eqref{10n}.

Тогда в \eqref{9n} матрица $J_2$ имеет на главной диагонали хотя бы один нуль.

Если  при этом матрицы $J_1$, $J_2$ содержат нуль на одной и той же строке,

то \mbox{$\det J'=\det J=0$}, что противоречит определению  матрицы $J$.

 Если же на некоторой строке матрица $J_2$ имеет нуль, а $J_1$

 "--- ненулевой элемент, то получаем противоречие \eqref{10n}.

 Следовательно,  $\det B\ne 0$. $\triangleleft$



\smallskip



%{\it Ниже приведена схему доказательства  пункта} 2). Из нее будет вытекать

%справедливость пунктов 1) и 3).

Сделаем в \eqref{7n} подстановки

$A=QJ_1Q^{-1},$ $B=QJ_2Q^{-1},$

тогда \eqref{7n} примет следующий вид:

\begin{multline}

\left[\,QJ_2^2Q^{-1}+QJ_2Q^{-1}QJ_1Q^{-1}QJ_2^{-1}Q^{-1}QJ_1Q^{-1}\,\right]

\,\frac{\partial^2 u}{\partial x^2}-

\\

-\left[\,QJ_2Q^{-1}QJ_1Q^{-1}QJ_2^{-1}Q^{-1}+QJ_1Q^{-1}\,\right]

\,\frac{\partial^2 u}{\partial x\partial y}

 +\frac{\partial^2 u}{\partial y^2}=

\\

 =Q\left[\,J_2^2+J_1'\,\right]Q^{-1}\,\frac{\partial^2 u}{\partial x^2}-

2QJ_1''Q^{-1}\,\frac{\partial^2 u}{\partial x\partial y}

 +\frac{\partial^2 u}{\partial y^2}=0.

 \label{12n}

 \end{multline}



В \eqref{12n} символами

$$

J_1'=J_2J_1J_2^{-1}J_1,\quad

J_1''=\frac{1}{2}\,\bigr(J_2J_1J_2^{-1}+J_1\bigr)

$$

 обозначены нижне-треугольные матрицы,  которые  на соответствующих

строках главной диагонали имеют {\it те же элементы}, {\it что и матрицы}

$J_1^2$, $J_1$ соответственно (но в общем случае не равны им).



Сделаем в нижнюю строку

 \eqref{12n} подстановку

 \begin{equation}

 u=Q\cdot p=Q\cdot(p_1,\ldots,p_n)^\top,

 \label{13n}

 \end{equation}

 где $p_k(z),$ $k=1,\ldots,n $ "---   скалярные комплексные функции. 

 Затем    умножим \eqref{12n} слева на $Q^{-1}$:

\begin{equation}

\bigr(J_2^2+J_1'\bigr)\,\frac{\partial^2 p}{\partial x^2}-

2J_1''\,\frac{\partial^2 p}{\partial x\partial y}+

\frac{\partial^2 p}{\partial y^2}=0.

\label{14n}

\end{equation}





Пусть $\xi_{1k},$ $\xi_{2k}\in\mathbb{C}$

"--- собственные числа матриц  $A$ и $B$

соответственно, стоящие на $k$-той строке  матриц $J_1$, $J_2$ в~\eqref{9n}. 

Пусть также  $L_k(\cdot)$ "--- линейные функции своих переменных.

Тогда систему  \eqref{14n} можно записать в следующем виде:

\begin{equation}

\begin{array}{l}

\displaystyle 

 \bigr(\xi_{11}^2+\xi_{21}^2\bigr)\,\frac{\partial^2 p_1}{\partial x^2}-

2\xi_{11}\,\frac{\partial^2 p_1}{\partial x\partial y}+

\frac{\partial^2 p_1}{\partial y^2}=0,\quad k=1;

\\

\displaystyle 

 \bigr(\xi_{1k}^2+\xi_{2k}^2\bigr)\,\frac{\partial^2 p_k}{\partial x^2}-

2\xi_{1k}\,\frac{\partial^2 p_k}{\partial x\partial y}+

\frac{\partial^2 p_k}{\partial y^2}=

\\

\displaystyle 

\hspace{1.4cm}

=L_k\Bigl(\frac{\partial^2 p_1}{\partial x^2}\,,\ldots,

\frac{\partial^2 p_{\scriptstyle k-1}}{\partial x^2},

\frac{\partial^2 p_1}{\partial x\partial y}\,,\ldots,

\frac{\partial^2 p_{\scriptstyle k-1}}{\partial x\partial y}\Bigr),

\quad k=2,\ldots,n.

\end{array}

\label{15n}

\end{equation}



Для {\it левой части}

 каждого из уравнений \eqref{15n} составим характеристический полином

(см. \cite{16Nk-Nk1})

$$

\lambda^2-2\xi_{1k}\lambda+\bigr(\xi_{1k}^2+\xi_{2k}^2\bigr)=0,

\quad k=1,\ldots,n,

$$

и найдем его корни:

\begin{equation}

\lambda_{1k}=\xi_{1k}+i\,\xi_{2k},\quad \lambda_{2k}=\xi_{1k}-i\,\xi_{2k},\quad

\xi_{1k},\,\xi_{2k}\in\mathbb{C},\quad k=1,\ldots,n.

\label{16n}

\end{equation}





В силу \eqref{9n}  $\lambda_{1k}$  в \eqref{16n} будет собственным числом

матрицы $J,$ а $\lambda_{2k}$ "--- собственным  числом матрицы $\overline{J},$

причем оба этих числа   стоят в $k$-той строке матриц $J'$, $J''$

соответственно.



Как показано в \cite{16Nk-Nk1}, если правая часть  уравнения вида $\eqref{15n}$

{\it равна нулю}  и при этом $\lambda_{1k}\ne\lambda_{2k}$ в \eqref{16n},

 то   общее решение такого  однородного

  уравнения   с учетом обозначений

 \eqref{2n} и \eqref{16n}   имеет вид

 \begin{equation}

p_k(z)=f_{\lambda_{\scriptstyle 1k}}(z)+f_{\lambda_{\scriptstyle 2k}}(z),

\quad k=1,\ldots,n,

\label{17n}

\end{equation}

то есть  является суммой произвольных

$\lambda_{1k}$- и $\lambda_{2k}$-голоморфных функций (см. определение~\hyperlink{nik:def2}{2}).



 %С учетом формулы  \eqref{16n}

%функции $L_k(\cdot)$ в \eqref{14n} будут

%линейными выражениями от вторых производных

%$\lambda$-голоморфных функций

%$f_{\lambda_{\scriptstyle 1l}},f_{\lambda_{\scriptstyle 2l}},\,$ $l=1,\ldots,k-1,$ $k\ge 2.$

% {\it По условию пункта} 2)

% среди чисел $\lambda_{1k},\lambda_{2k},$ $k=1,\ldots,n$ {\it нет одинаковых}.

% Поэтому в силу той же формулы \eqref{16n} для общего решения,

%   после подстановки функции

% $f_{\lambda_{\scriptstyle 1l}}$  или $f_{\lambda_{\scriptstyle 2l}}$

% при $l\le k-1$

%  в левую часть $k$-ой строки  \eqref{14n}

%  справа получим сумму ее вторых производных

%  с ненулевыми коэффициентами.

%  Но в силу  \eqref{2n} вторые производные

%  $\lambda$-голоморфной функции

%  по разным переменным отличаются одна от другой на константу.

% Таким образом,

%    общие решения {\it неоднородных } уравнений \eqref{14n}

%при $k\ge 2$ можно найти в виде

%\begin{equation}

%p_k=f_{\lambda_{\scriptstyle 1k}}+f_{\lambda_{\scriptstyle 2k}}+

%\widetilde{L}_k\left(f_{\lambda_{\scriptstyle 11}},f_{\lambda_{\scriptstyle 21}},\ldots,

%f_{\lambda_{\scriptstyle 1k-1}},f_{\lambda_{\scriptstyle 2k-1}}\right),\quad

%k=2,\ldots,n,

%\label{17n}

%\end{equation}

%где $\widetilde{L}_k(\cdot)$ --- линейные функции своих переменных.



С учетом подстановки    \eqref{13n}

граничное условие \eqref{3n} можно записать в следующем виде:

 \begin{multline}

\bigr(p_1(z),\ldots,p_n(z)\bigr)^\top\bigr|_\Gamma=

Q^{-1}u(x,y)\bigr|_\Gamma=\\

=Q^{-1}\cdot\bigr[\mbox{Re}\,\phi(z)\bigr]\bigr|_\Gamma=

Q^{-1} \cdot\bigr(\varphi_1,\ldots,\varphi_n\bigr)^\top.

\rule{0pt}{5mm}

\label{18n}

\end{multline}





{\it Пусть  выполнены условия пункта} \hyperlink{nik:th4-2}{2)} {\it  и неравенства $\eqref{10n}.$ }

Тогда правая часть \eqref{18n} будет нулевой.



Допустим, что существует нетривиальное решение $\phi_0=u_0+iv_0$

однородной задачи Шварца.

Тогда в силу \eqref{13n} функция $p=Q^{-1}\cdot u_0$ тоже непостоянна.

Заметим, что с учетом теоремы \hyperlink{nik:th3}{3} и формулы \eqref{17n} для общего решения функция $p_1(z)$ равна константе. 

Отсюда в~\eqref{15n} функция $L_2({}\cdot{})\equiv 0,$ то есть в силу \eqref{17n} и теоремы~\hyperlink{nik:th3}{3}  $p_2(z)\equiv\mbox{const}.$

Применяя далее тривиальную индукцию, получаем, что функция $p(z)$ постоянна, что противоречит сделанному предположению. Это противоречие и {\it доказывает пункт {\rm \hyperlink{nik:th4-2}{2)}} настоящей теоремы. }



\smallskip



 {\it Пусть теперь  выполнены условия пункта {\rm \hyperlink{nik:th4-1}{1)}} и неравенства $\eqref{10n}.$ }

В силу диагональности жордановых форм матриц $A,B$ система \eqref{15n}

будет однородной, то есть все функции $L_k(\cdot)$  равны  нулю.

Поэтому с  учетом формул

 \eqref{10n}, \eqref{17n}, \eqref{18n}

 и теоремы~\hyperlink{nik:th2}{2} можно однозначно найти две  функции $f_{\lambda_{\scriptstyle 1k}}(z)$, $f_{\lambda_{\scriptstyle 2k}}(z)

\in H^\sigma(\overline{D}),$  то есть $p_k(z),$ $k=1,\ldots,n.$





В результате найдена

вещественная функция

$u=Q\cdot (p_1,\ldots,p_n)\in H^\sigma(\overline{D}),$

которая   есть не что иное, как решение {\it задачи Дирихле}

\eqref{8n} для вещественной

системы \eqref{7n}

 с заданной граничной функцией $\varphi(z)\in H^\sigma(\Gamma).$



\smallskip



    Для доказательства  существования решения {\it задачи Шварца }

нужно  восстановить функцию $v=v(x,y)$ по уже известной $u(x,y),$

  используя для определенности первую  из формул \eqref{6n}. 

  Тогда   функция $\phi(z)=u(x,y)+iv(x,y)$ будет решением задачи Шварца

  \eqref{3n} с той же граничной функцией $\varphi(z),$ причем $\phi(z)$ будет соответствовать данной матрице $J=A+Bi.$





С  учетом формул \eqref{6n}, \eqref{13n},  \eqref{17n} и \eqref{2n}  имеем

\begin{multline}

\frac{\partial v}{\partial x}=B^{-1}A\,

\frac{\partial u}{\partial x}

-B^{-1}\,\frac{\partial u}{\partial y}=

B^{-1}AQ\cdot\frac{\partial p}{\partial x}

-B^{-1}Q\cdot\frac{\partial p}{\partial y}=

\\

 = B^{-1}AQ\cdot\frac{\partial p}{\partial x}-B^{-1}Q\cdot

\begin{pmatrix}

{\displaystyle \frac{\partial f_{\lambda_{\scriptstyle 11}}}{\partial y}+

\frac{\partial f_{\lambda_{\scriptstyle 21}}}{\partial y} }\\

\cdots\cdots\cdots\cdots\cdots\\

{\displaystyle\frac{\partial f_{\lambda_{\scriptstyle 1n}}}{\partial y}+

\frac{\partial f_{\lambda_{\scriptstyle 2n}}}{\partial y}\rule{0pt}{5mm} }

\end{pmatrix}=

\\

 =B^{-1}AQ\cdot

\begin{pmatrix}

{\displaystyle \frac{\partial f_{\lambda_{\scriptstyle 11}}}{\partial x}+

\frac{\partial f_{\lambda_{\scriptstyle 21}}}{\partial x} }\\

\cdots\cdots\cdots\cdots\cdots\\

{\displaystyle\frac{\partial f_{\lambda_{\scriptstyle 1n}}}{\partial x}+

\frac{\partial f_{\lambda_{\scriptstyle 2n}}}{\partial x} \rule{0pt}{5mm} }

\end{pmatrix}-

B^{-1}Q\cdot

\begin{pmatrix}

{\displaystyle\lambda_{11}

\frac{\partial f_{\lambda_{\scriptstyle 11}}}{\partial x}+

\lambda_{21}\frac{\partial f_{\lambda_{\scriptstyle 21}}}{\partial x} }\\

 \cdots\cdots\cdots\cdots\cdots\cdots\cdots\\

    {\displaystyle\lambda_{1n}

\frac{\partial f_{\lambda_{\scriptstyle 1n}}}{\partial x}+

\lambda_{2n}\frac{\partial f_{\lambda_{\scriptstyle 2 n}}}{\partial x}

\rule{0pt}{5mm}  }

\end{pmatrix},

\label{19n}

\end{multline}

где в последних  скобках производная по $y$ замена на производную по $x$

 с~учетом определения~\hyperlink{nik:def2}{2}  $\lambda$-голоморфной функции.



 Левая и правая части  \eqref{19n}

  представляют собой производные некоторых

  функций  по одной и той же переменной  $x.$ Следовательно, эти функции

  совпадают   с точностью до

   вектор-функция $r=r(y),$ которая  зависит только от  переменной $y$:

\begin{equation}

  v(x,y)= B^{-1}AQ\cdot

\begin{pmatrix}

f_{\lambda_{\scriptstyle 11}}+

f_{\lambda_{\scriptstyle 21}}\\

\cdots\cdots\cdots\\

f_{\lambda_{\scriptstyle 1n}}+

f_{\lambda_{\scriptstyle 2n}}

\end{pmatrix}-

B^{-1}Q\cdot

\begin{pmatrix}

\lambda_{11}

f_{\lambda_{\scriptstyle 11}}+

\lambda_{21}f_{\lambda_{\scriptstyle 21}}\\

     \cdots\cdots\cdots\cdots\cdots\cdots\\

\lambda_{1n}

 f_{\lambda_{\scriptstyle 1n}}+

\lambda_{2n} f_{\lambda_{\scriptstyle 2 n}}

\end{pmatrix}+r(y).

\label{20n}

\end{equation}



Функция $\phi=u+iv,$ где $u,v$ имеют вид \eqref{13n} и \eqref{20n}

 соответственно, по

 построению будет решением уравнения

 \eqref{1n} при некотором подборе  $r(y).$

 Функция $r(y)$ определяется с помощью второго

уравнения в \eqref{6n}. Чтобы его не использовать,

покажем, что  $r(y)$ в \eqref{20n} {\it постоянна.}



\smallskip



$\triangleright$

Действительно,

 $r(y)\in C^1(D),$ так как дифференцируемы функции $\phi$ и $\phi-ri.$

Непосредственная подстановка показывает, что $\phi-ri$ удовлетворяет \eqref{1n},

то есть и функция $r(y)$ будет аналитической по Дуглису, откуда вытекает, что

 $r(y)\equiv\rm{const}.$ $\triangleleft$



\smallskip



{\it В итоге  формулы  $\eqref{13n}$ и $\eqref{20n}$ определяют  решение

$\phi=u+iv$ неоднородной задачи Шварца } \eqref{3n}.



Поскольку все функции, входящие  в \eqref{20n}, а  также функция $u(x,y),$

непрерывны по Гельдеру в

$\overline{D},$ то и $\phi(z)\in H^\sigma(\overline{D}).$

Разумеется, те же результаты будут иметь место, если

использовать вторую  из формул \eqref{6n}.



 {\it Единственность} полученного решения доказывается

 аналогично  пункту~\hyperlink{nik:th4-2}{2)}.

\end{proof}





\smallskip



\Section{Построение  матриц $\boldsymbol J$, для которых

справедлива теорема~\hyperlink{nik:th4}{4}}

Обозначим

\begin{equation}

J=P_1(A)+P_2(A)i,\quad A\in\mathbb{R}^{n\times n},

\label{21n}

\end{equation}

где $P_1$, $P_2$ "---  матричные полиномы с вещественными коэффициентами.

 Очевидно, что

 матрицы $A$, $P_1(A)$, $P_2(A)$, $J$  в~\eqref{21n} имеют общий базис $Q,$ приводящий их к  треугольному виду.

В качестве $Q$ можно взять жорданов базис матрицы $A.$



Для   матриц  \eqref{21n}  нужно проверить

выполнение  критерия \eqref{10n} или  \eqref{11n}.

Какой именно из двух критериев нужно проверять, 

зависит от формы записи матрицы $J.$ 

Если $J=(A+Bi)^m,$ $m=2, 3, \ldots,$ то удобнее применять \eqref{10n}, а если $J=A+Bi,$

то более удобен  критерий \eqref{11n}.



Всюду в этом пункте будем обозначать через

 $\xi_k+i\delta_k,$ $\xi_k$, $\delta_k\in\mathbb{R},$

 \mbox{$k=1, \ldots, n,$} собственные числа вещественной  матрицы    $A.$

 Комплексные числа $\lambda_{1k}$, $\lambda_{2k}$ определены перед формулой

 \eqref{10n},

а $\lambda_k$, $\mu_k$ "--- перед формулой~\eqref{11n}.



\smallskip



$1^\circ.$ Для матрицы

 $J^*=\alpha E+Ai,$ $\alpha\in\mathbb{R}$ все числа  $\xi_k\ne 0,$

так как в~противном случае матрица $J^*$ будет

 иметь вещественные собственные числа,

что противоречит определению~\hyperlink{nik:def1}{1}. 

Поэтому

$$

|\mathop{\rm Im} \lambda_k|=0<|\mathop{\rm Re} \mu_k|=|\xi_k|\ne 0,\quad k=1, \ldots, n,

$$

то  есть  выполнен критерий \eqref{11n}.  Таким образом, имеет место



\smallskip



\hypertarget{nik:lemma2}{}



\begin{lemma}[2]Для функций$,$ $J$-аналитических с матрицами  $J^*=\alpha E+Ai,$ $\alpha\in\mathbb{R},$

      теорема {\rm \hyperlink{nik:th4}{4}} справедлива при любых   значениях $\alpha\in\mathbb{R}.$

\end{lemma}



\smallskip



Отметим, что утверждение леммы~\hyperlink{nik:lemma2}{2}  было доказано   в \cite{2Nk-Nk2}.

 Однако в \cite{2Nk-Nk2} более общие случаи не рассматривались  и~критерий \eqref{10n}, \eqref{11n} получен не был.



\smallskip



$2^\circ.$ Рассмотрим матрицу $\widetilde{J}=A+\alpha iE,$

$\alpha\in\mathbb{R}.$ Пусть выполнено неравенство

\begin{equation}

\max\limits_{1\le k \le n}|\delta_k|<|\alpha|.

\label{22n}

\end{equation}

Тогда

 $$

|\mathop{\rm Im}\lambda_k|=|\delta_k|<|\alpha|=

|\mathop{\rm Re}\mu_k|, \quad k=1,\ldots,n,

$$

то есть имеет место формула \eqref{11n}.

Следовательно,  справедлива



\smallskip



 \begin{lemma}[3] Для функций$,$ $J$-аналитических с матрицами

   $\widetilde{J}=A+\alpha iE,$ $\alpha\in\mathbb{R},$

  справедлива теорема  {\rm \hyperlink{nik:th4}{4},} 

  если верно неравенство $\eqref{22n}.$

  \end{lemma}



\smallskip



  $3^\circ.$ Рассмотрим матрицу

  $J^{*2}=(\alpha E+Ai)^2.$ Здесь

  \begin{gather}

  \lambda_{1k}=\bigr(\alpha+(\xi_k+i\delta_k)i\bigr)^2=

  (\alpha-\delta_k+\xi_ki)^2,\notag\\

  \lambda_{2k}=\bigr(\alpha-(\xi_k+i\delta_k)i\bigr)^2=

  (\alpha+\delta_k-\xi_ki)^2.\notag

  \end{gather}







   Проверим выполнение критерия (\ref{10n}):

\begin{equation}

\left(\mathop{\rm Im} \lambda_{1k} \right)\cdot

\left(\mathop{\rm Im} \lambda_{2k} \right)%=\\

=-2\xi_k(\alpha-\delta_k)\cdot

2\xi_k(\alpha+\delta_k)=-4\xi_k^2\bigr(\alpha^2-\delta_k^2\bigr)<0,

  \end{equation}

если $\xi_k\ne 0 $ и  верно \eqref{22n}.

В итоге доказана



\smallskip



\begin{lemma}[4]Для функций$,$ $J$-аналитических с матрицами

 $$

 J^{*2}=(\alpha E+Ai)^2,\quad

 \alpha\in\mathbb{R},\quad \xi_k\ne 0,

 $$

  справедлива   теорема {\rm \hyperlink{nik:th4}{4},} 

  если имеет место  неравенство $\eqref{22n}.$

  \end{lemma}



\smallskip



\Section{Матрицы $\boldsymbol J$, отличные от (\ref{21n}), для которых

выполнен критерий (\ref{11n})}

Пусть

 $J^*_1$, $J^*_2$ "--- блочно-диагональные матрицы

  размера $l\times l,$ блоками которых являются жордановы клетки

 (в общем случае комплексные). 

 Пусть

 $\widetilde{J}_1$, $\widetilde{J}_2$ "--- блочно-диагональные матрицы

 размера $(n-2l)\times(n-2l), $

  состоящие  из вещественных жордановых клеток.

   В обеих случаях жордановы клетки могут быть одномерными.

      Пусть матрица $Q^*\in\mathbb{C}^{n\times l},$

а $\widetilde{Q}$ "--- вещественная матрица размера $n\times (n-2l).$

     Образуем матрицу $J$ по следующему правилу:

        \begin{multline}

         J_1={\rm diag}\,\bigr(\,

 J^*_1,\,\overline{J^*_1},\,\widetilde{J}_1\,\bigr),\quad

J_2={\rm diag}\,\bigr(\,

 J^*_2,\,\overline{J^*_2},\,\widetilde{J}_2\,\bigr),\quad

   Q=\bigr(\,

 Q^*,\,\overline{Q^*},\,\widetilde{Q}\,\bigr),\\

  A=QJ_1Q^{-1},\quad B=QJ_2Q^{-1},\quad J=A+Bi.

   \label{23n}

\end{multline}





 Покажем, что  матрицы $A$, $B$ в \eqref{23n} "---  вещественные.



\smallskip



 $\triangleright$

 Действительно, $\overline{Q^{-1}}=\left(\overline{Q}\right)^{-1},$

 что вытекает из алгоритма построения обратной матрицы.

 Поэтому

 \begin{equation}

 \overline{A}=\overline{QJ_1Q^{-1}}=

\overline{Q}\cdot\overline{J_1}\cdot\overline{Q^{-1}}=

\overline{Q}\cdot\overline{J_1}\cdot\left(\overline{Q}\right)^{-1}.

 \label{24n}

\end{equation}



Как известно, жорданов базис $Q$ и жорданова  форма $J_1$ матрицы $A$

определены с точностью до перестановки жордановых клеток в $J_1$ и

соответствующих им столбцов матрицы $Q.$

Из \eqref{23n}  следует, что операция сопряжения матриц $Q$ и $J_1$

как раз и производит такую перестановку.  Следовательно,

$\overline{Q}$ и $\overline{J}_1$ в \eqref{24n}

 определяют ту же самую матрицу $A,$

то есть  $\overline{A}=A.$

  Аналогично доказывается равенство

  $\overline{B}=B.$ $\triangleleft$

\smallskip



 Таким образом,

формулы \eqref{23n} описывают альтернативный \eqref{21n}

способ построения  матриц $J$

из условий теоремы~\hyperlink{nik:th4}{4}.





Имеет место очевидное следствие из теоремы~\hyperlink{nik:th4}{4}, а именно



\smallskip



\hypertarget{nik:lemma5}{}



\begin{lemma}[5]Пусть $\lambda_k,\mu_k,$

 $k=1,\ldots,n$ --- собственные  числа

   матриц $J_1,J_2$  $\eqref{23n}$

  соответственно$,$  стоящие на $k$-той строке$.$ Пусть выполнен критерий

   $\eqref{11n}.$

  Тогда для  функций, $J$-аналитических с матрицами

  $J$ $\eqref{23n},$ справедлива теорема $4.$

  \end{lemma}



\smallskip



Очевидно, что применение критерия \eqref{10n}

для матриц вида \eqref{23n}  затруднительно.

В заключение приведем следующий пример, иллюстрирующий

лемму~\hyperlink{nik:lemma5}{5}.





\smallskip



 \begin{example}[2]Пусть числа

 $\lambda,\mu,\eta$ имеют ненулевые комплексные части;

 векторы $\bold x,\bold y\in\mathbb{C}^4.$

   Положим

\begin{equation}

J_1=

\begin{pmatrix}

\lambda&0&0&0\\

0&\lambda&0&0\\

0&0&\overline{\lambda}&0\\

0&0&0&\overline{\lambda}

\end{pmatrix},\quad

J_2=\begin{pmatrix}

\mu&0&0&0\\

0&\eta&0&0\\

0&0&\overline{\mu}&0\\

0&0&0&\overline{\eta}

\end{pmatrix},\quad

Q=\bigr(\bold x,\bold y,\overline{\bold x},\overline{\bold y}\bigr).

\label{25n}

\end{equation}

\end{example}



Матрицы   \eqref{25n}   имеют вид \eqref{23n},

причем $J_1$, $J_2$ "--- диагональны.

Поэтому для  функций, $J$-аналитических с матрицами

$$

J=A+Bi=QJ_1Q^{-1}+QJ_2Q^{-1}i,

$$

справедливы утверждения  леммы \hyperlink{nik:lemma5}{5} и теоремы \hyperlink{nik:th4}{4},

если

$$

 |\mathop{\rm Im} \lambda|<\min\bigr\{ |\mathop{\rm Re} \mu|,\,|\mathop{\rm Re} \eta|

 \bigr\},

 $$

то есть если  выполнен  критерий \eqref{11n}.



\smallskip



\Section[n]{Заключение} Критерий существования и единственности

решений задачи Шварца

 \eqref{10n}, \eqref{11n}

довольно прост для проверки и может быть

 применен к различным матрицам вида  \eqref{21n} и \eqref{23n}.

  Поскольку задача Шварца, как было отмечено выше,

   равносильна задаче Дирихле для

  многих  типов эллиптических

   систем второго порядка в частных производных,

   то неравенства \eqref{10n}, \eqref{11n} применимы и для

изучения   этой задачи.







\thanks{{\bf Декларация о финансовых и других взаимоотношениях.}

Работа выполнена при поддержке Министерства образования и~науки Российской Федерации в~рамках государственного задания (проект 1.857.2014/K).

Автор несет полную ответственность за предоставление окончательной версии рукописи в печать.

Окончательная версия рукописи была одобрена автором.

Автор не получал гонорар за статью.

}











\thanks{{\bf ORCID}



{\bf Владимир Геннадьевич Николаев:} \url{http://orcid.org/0000-0003-0274-5826}

}



\RusRef



%% Благодарности, гранты и прочее



\begin{thebibliography}{99}



\RBibitem{1Nk-S1}

\paper Задача Шварца для функций, аналитических по Дуглису

\by Солдатов~А.~П.

\jour Совр. математика и ее приложения

\yr 2010

\vol 67

\volinfo Уравнения с частными производными

\pages 99--102

\zmath{https://zbmath.org/?q=an:1220.30066}



\RBibitem{2Nk-Nk2}

\paper О  решении задачи Шварца для $J$-аналитических функций в областях,

ограниченных контуром Ляпунова

\by   Николаев~В.~Г., Солдатов~А.~П.

\jour Диффер.  уравн.

\yr 2015

\issue 7

\pages 965--969

\zmath{https://zbmath.org/?q=an:06493795}



\RBibitem{4Nk-S3}

\paper Интегральное представление функций, аналитических по Дуглису

\by Солдатов~А.~П.

\jour Вестн. СамГУ. Естественнонаучн. сер.

\yr 2008

\issue 8/1(67)

\pages 225--234









\RBibitem{3Nk-S2}

\paper Гипераналитические функции и их приложения

\by Солдатов~А.~П.

\jour Совр. математика и ее приложения

\yr 2004

\vol 15

\volinfo Теория функций

\pages 142--199







\RBibitem{5Nk-S5}

\paper Пространство Харди решений эллиптических систем первого порядка

\by Солдатов~А.~П.

\jour Докл. РАН

\yr 2007

\vol 416

\issue 1

\pages 26--30

\elib{http://elibrary.ru/item.asp?id=9533780}





\RBibitem{6Nk-S6}

\paper Эллиптические системы высокого порядка

\by Солдатов~А.~П.

\jour Диффер. уравн.

\yr 1989

\vol 25

\issue 1

\pages 136--142

\zmath{https://zbmath.org/?q=an:0697.35025}





\RBibitem{7Nk-B1}

\by Бицадзе~А.~В.

\paper О~единственности решения задачи Дирихле для эллиптических

уравнений с~частными производными

\jour УМН

\yr 1948

\vol 3

\issue 6(28)

\pages 211--212

\mathnet{http://mi.mathnet.ru/umn8780}

\mathscinet{http://www.ams.org/mathscinet-getitem?mr=27409}

\zmath{https://zbmath.org/?q=an:0041.21703}





\RBibitem{8Nk-B2}

\by Бицадзе~А.~В.

\book Основы теории аналитических функций комплексного переменного

\publaddr М.

\publ Наука

\yr 1972

\totalpages 347

\zmath{https://zbmath.org/?q=an:0259.30001}



\RBibitem{9Nk-B3}

\by Бицадзе~А.~В.

\book Краевые задачи для эллиптических уравнений второго порядка

\publaddr М.

\publ Наука

\yr 1966

\totalpages 298

\zmath{https://zbmath.org/?q=an:0148.08904}





\RBibitem{10Nk-Bo-1}

\paper Теория обобщенного аналитического вектора

\by Боярский~Б.~В.

\jour Annales Polonici Mathematici

\yr 1965

\issue 3

\vol 17

\pages 281--320

\elink \url{https://eudml.org/doc/265098}





\RBibitem{11Nk-V1}

\by Векуа~И.~Н.

\book Обобщенные аналитические функции

\publaddr М.

\publ Наука

\yr 1988

\totalpages 328

\zmath{https://zbmath.org/?q=an:0698.47036}





\RBibitem{12Nk-J1}

\paper Краевые задачи типа Бицадзе--Самарского для эллиптических    в смысле Дуглиса--Ниренберга систем

\by  Жура~Н.~А.

\jour Диффер. уравн.

\yr 1992

\vol 28

\issue 1

\pages 81--91

\zmath{https://zbmath.org/?q=an:0780.35031|an:0788.35037}





\RBibitem{13Nk-J2}

\paper Общая краевая задача для эллиптических в смысле Дуглиса--Ниренберга систем в областях с гладкой границей

\by  Жура~Н.~А.

\jour Изв. РАН

\yr 1994

\issue 58

\issue 1

\pages 22--44

\zmath{https://zbmath.org/?q=an:0831.35050}



\RBibitem{14Nk-Mus}

\by Мусхелишвили~Н.~И.

\book Сингулярные интегральные уравнения. Граничные задачи теории функций и некоторые их приложения к математической физике

\publaddr М.

\publ Наука

\yr 1968

\totalpages 511

\zmath{https://zbmath.org/?q=an:03277892}



\RBibitem{15Nk-Nk2}

\by Николаев~В.~Г.

\paper О некоторых  свойствах $J$-аналитических функций

\jour Вестн. СамГУ. Естественнонаучн. сер.

\yr 2013

\issue 3(104)

\pages 25--32

\mathnet{http://mi.mathnet.ru/vsgu323}



\RBibitem{16Nk-Nk1}

\paper Результаты о совпадении $\lambda$-   и $\mu$-голоморфных функций на границе области и    их приложения к краевым задачам

\by   Николаев~В.~Г., Панов~Е.~Ю.

\inbook Проблемы математического  анализа

\bookinfo Межвузовский международный сборник 

\eds Н.~Н.~Уральцева

\yr 2013

\volinfo Вып.~74

\pages 123--132

\publaddr Новосибирск 

\publ Тамара Рожковская





\end{thebibliography}



\RuDate



\newpage





\makeentitle



\thanks{{\bf Declaration of Financial and Other Relationships.}

This work was supported by the Ministry of education and science of the Russian Federation within the framework of the base part of the state task  (Project 1.857.2014/K).

The author is absolutely responsible for submitting the final manuscript in print.

The author has approved the final version of manuscript.

The author has not received any fee for the article.

}













\thanks{{\bf ORCID}



{\bf  Vladimir G. Nikolaev:} \url{http://orcid.org/0000-0003-0274-5826}

}





\pagebreak



\EngRef





\begin{thebibliography}{99}



\Bibitem{1Nk-S1-en}

\by Soldatov~A.~P.

\paper The Schwarz problem for Douglis analytic functions

\jour J. Math. Sci.

\vol 173

\issue 2

\pages 221--224 

\yr 2011

\crossref{10.1007/s10958-011-0244-7}

\zmath{https://zbmath.org/?q=an:1220.30066}







\Bibitem{2Nk-Nk2-en}

\by Nikolaev~V.~G., Soldatov~A.~P.

\paper On the solution of the Schwarz problem for $J$-analytic functions in a domain bounded by a Lyapunov contour 

\jour Differ. Equ. 

\vol 51

\issue 7

\pages 962--966 

\yr 2015

\zmath{https://zbmath.org/?q=an:06493795}

\crossref{10.1134/S0012266115070150}



\Bibitem{4Nk-S3-en}

\lang In Russian

\by Soldatov~A.~P.

\paper Integral representation of Douglis analytic functions

\jour Vestnik SamGU. Estestvenno-Nauchnaya Ser.

\yr 2008

\issue 8/1(67)

\pages 225--234





\Bibitem{3Nk-S2-en}

\lang In Russian

\paper Hyperanalytic functions and their applications

\by Soldatov~A.~P.

\jour Sovr. Mat. Pril.

\yr 2004

\vol 15

\pages 142--199





\Bibitem{5Nk-S5-en}

\paper The Hardy space of solutions to first-order elliptic systems

\by Soldatov~A.~P.

\jour Doklady Mathematics

\yr 2007

\vol 76

\issue 2

\pages 660--664

\elib{http://elibrary.ru/item.asp?id=13536051}

\crossref{10.1134/S1064562407050067}





\Bibitem{6Nk-S6-en}

\by Soldatov~A.~P.

\paper High-order elliptic systems

\jour Differ. Equ.

\vol 25

\issue 1

\pages 109--115 

\yr 1989

\zmath{https://zbmath.org/?q=an:0697.35025}





\Bibitem{7Nk-B1-en}

\lang In Russian

\by Bitsadze~A.~V.

\paper On the uniqueness of the solution of the Dirichlet problem for elliptic partial differential equations

\jour Uspekhi Mat. Nauk

\yr 1948

\vol 3

\issue 6(28)

\pages 211--212



\Bibitem{8Nk-B2:ge}

\lang In German

\by Bizadse~A.~W.

\book Grundlagen der Theorie analytischer Funktionen

\serial Mathematische Lehrb\"ucher und Monographien

\bookinfo I.~Abteilung, Band 23

\publaddr Berlin

\publ Akademie-Verlag

\totalpages 186 

\yr 1973

\zmath{https://zbmath.org/?q=an:0259.30002}



\Bibitem{9Nk-B3:en}

\by Bitsadze~A.~V.

\book Boundary value problems for second order elliptic equations

\serial North-Holland Series in Applied Mathematics and Mechanics

\vol 5

\publaddr Amsterdam

\publ North-Holland Publ. Company

\totalpages 211 

\yr 1968

\zmath{https://zbmath.org/?q=an:0167.09401}









\Bibitem{10Nk-Bo-1:en}

\lang In Russian

\paper The theory of generalized analytic vector

\by Boyarskii~B.~V.

\jour Annales Polonici Mathematici

\yr 1965

\issue 3

\vol 17

\pages 281--320

\elink \url{https://eudml.org/doc/265098}









\Bibitem{11Nk-V1:en}

\by Vekua~I.~N.

\book Generalized analytic functions

\zmath{https://zbmath.org/?q=an:0100.07603}

\serial International Series of Monographs on Pure and Applied Mathematics

\vol 25

\publaddr Oxford 

\publ Pergamon Press 

\yr 1962

\totalpages xxvi+668 







\Bibitem{12Nk-J1:en}

\by Zhura~N.~A.

\paper Bitsadze-Samarskij type boundary-value problems for Douglis-Nirenberg elliptical systems

\jour Differ. Equ.

\vol 28

\issue 1

\pages 79--88 

\yr 1992

\zmath{https://zbmath.org/?q=an:0780.35031|an:0788.35037}





\Bibitem{13Nk-J2:en}

\by Zhura~N.~A.

\paper A general boundary value problem for systems elliptic in the Douglis- Nirenberg sense in domains with smooth boundary

\jour Russ. Acad. Sci., Izv., Math. 

\vol 44

\issue 1

\pages 21--42 

\yr 1995

\zmath{https://zbmath.org/?q=an:0831.35050}

\crossref{10.1070/IM1995v044n01ABEH001584}



\Bibitem{14Nk-Mus:en}

\by Muskhelishvili~N.~I.

\book Singular integral equations

\zmath{https://zbmath.org/?q=an:03277891}

\publaddr Groningen

\publ Wolters-Noordhoff Publ.

\totalpages 447 

\yr 1967



\Bibitem{15Nk-Nk2:en}

\lang In Russian

\by   Nikolaev~V.~G.

\paper On some properties of $J$-analytical functions

\jour Vestnik SamGU. Estestvenno-Nauchnaya Ser.

\yr 2013

\issue 3(104)

\pages 25--32



\Bibitem{16Nk-Nk1:en}

\lang In Russian

\paper On coincidence of $\lambda$- and $\mu$-holomorphic functions on the boundary of a domain and applications to elliptic boundary value problems

\by   Nikolaev~V.~G., Panov~E.~Yu.

\inbook Problemy matematicheskogo  analiza \rm [Problems of mathematical analysis]

\eds N.~N.~Ural'tseva

\yr 2013

\volinfo Issue~74

\pages 123--132

\publaddr Novosibirsk 

\publ Tamara Rozhkovskaia







\end{thebibliography}



\EngDate

\Russian

\newpage





%\end{document}

%

